\documentclass[reqno]{amsart}
\usepackage[T1]{fontenc}
\usepackage{amsmath,amssymb,amsthm}
\usepackage{xr-hyper}
\usepackage[hidelinks]{hyperref}
\externaldocument[zr-]{20261125T001000--zeros-to-ramp-sums__workbook}
\externaldocument[gf-]{20261125T005000--galois-fourier-and-averaging__workbook}
\externaldocument[sb-]{20261125T020200--slopes-frames-and-bessel-poincare-series__workbook}
\externaldocument[co-]{20261126T020140--the-colength-of-a-hecke-image__workbook}
\externaldocument[gl-]{20261126T020150--gain-and-loss-at-an-inert-prime__workbook}
\externaldocument[as-]{20261125T021500--assembly-at-one-scale__workbook}

% Chapter number; \ref{chap} gives it.
\newcounter{chap}
% Chapter numbers, generated from 20261125T012000--grh-workbook__guide.org
% by scripts/make-chapters.py.  Do not edit: change the order of the
% index and rerun the script (./preview-site.sh and CI do this).
%
% Usage in a sheet, after \newcounter{chap}:
%   % Chapter numbers, generated from 20261125T012000--grh-workbook__guide.org
% by scripts/make-chapters.py.  Do not edit: change the order of the
% index and rerun the script (./preview-site.sh and CI do this).
%
% Usage in a sheet, after \newcounter{chap}:
%   % Chapter numbers, generated from 20261125T012000--grh-workbook__guide.org
% by scripts/make-chapters.py.  Do not edit: change the order of the
% index and rerun the script (./preview-site.sh and CI do this).
%
% Usage in a sheet, after \newcounter{chap}:
%   \input{chapters}\chapnumber{<sheet ID>}
\makeatletter
\providecommand\chapters@set[2]{\@namedef{chapters@#1}{#2}}
\chapters@set{20261125T001000}{1}% zeros-to-ramp-sums
\chapters@set{20261125T001100}{2}% inert-primes-and-estimate-b
\chapters@set{20261125T005000}{3}% galois-fourier-and-averaging
\chapters@set{20261125T005500}{4}% columns-from-a-sphere
\chapters@set{20261125T000000}{5}% slope-method
\chapters@set{20261126T030100}{6}% a-twist-for-each-prime
\chapters@set{20261125T003000}{7}% the-chain-and-the-square-root
\chapters@set{20261125T010000}{8}% rank-lemma-and-zero-estimate
\chapters@set{20261125T010500}{9}% poincare-series-and-near-orthogonality
\chapters@set{20261124T223000}{10}% elliptic-curves-for-the-proof
\chapters@set{20261125T009500}{11}% fock-model-and-theta-intertwiners
\chapters@set{20261124T223100}{12}% complex-multiplication-and-inert-reduction
\chapters@set{20261125T008000}{13}% one-cm-determinant
\chapters@set{20261124T214500}{14}% supersingular-theta-divisibility
\chapters@set{20261124T214600}{15}% the-isotropic-plane
\chapters@set{20261125T013700}{16}% pairings-on-a-formal-group-of-height-two
\chapters@set{20261125T013500}{17}% hecke-nondegeneracy-on-the-tree
\chapters@set{20261126T030200}{18}% the-bound-is-attained
\chapters@set{20261125T009000}{19}% kronecker-limit-and-chowla-selberg
\chapters@set{20261125T002000}{20}% definite-theta-lifts
\chapters@set{20261125T015500}{21}% siegel-weil-in-the-smallest-case
\chapters@set{20261125T002100}{22}% definite-theta-lifts-and-norms
\chapters@set{20261125T016500}{23}% hecke-operators-and-theta-lifts
\chapters@set{20261125T016700}{24}% divisibility-along-the-chain
\chapters@set{20261125T004000}{25}% denominators-and-sha
\chapters@set{20261125T004100}{26}% tate-twists-and-sha
\chapters@set{20261126T020430}{27}% the-minus-part-by-gauss-sums
\chapters@set{20261125T014100}{28}% congruence-modules
\chapters@set{20261125T014500}{29}% eisenstein-congruences-and-selmer-classes
\chapters@set{20261125T014700}{30}% pseudocharacters-and-the-length-bound
\chapters@set{20261125T011000}{31}% galois-representations-black-box
\chapters@set{20261126T030300}{32}% the-data
\chapters@set{20261126T020210}{33}% bessel-poincare-series-and-their-gram-operators
\chapters@set{20261126T030400}{34}% the-seesaw-identity
\chapters@set{20261126T030500}{35}% divisibility-of-the-composite-lift
\chapters@set{20261126T030600}{36}% theorem-a-short-forms
\chapters@set{20261126T030700}{37}% the-determinant-of-a-chain
\chapters@set{20261126T030800}{38}% the-cost-of-one-lift
\chapters@set{20261126T030900}{39}% denominators-of-projections-to-lifts
\chapters@set{20261126T031000}{40}% twists-with-few-congruences
\chapters@set{20261126T031100}{41}% theorem-c-and-the-main-inequality
\providecommand\chapnumber[1]{%
  \@ifundefined{chapters@#1}%
    {\PackageError{chapters}{Sheet #1 is not listed in the index}%
      {List it in 20261125T012000--grh-workbook__guide.org\MessageBreak
       and run scripts/make-chapters.py.}}%
    {\edef\thechap{\@nameuse{chapters@#1}}}}
\makeatother
\chapnumber{<sheet ID>}
\makeatletter
\providecommand\chapters@set[2]{\@namedef{chapters@#1}{#2}}
\chapters@set{20261125T001000}{1}% zeros-to-ramp-sums
\chapters@set{20261125T001100}{2}% inert-primes-and-estimate-b
\chapters@set{20261125T005000}{3}% galois-fourier-and-averaging
\chapters@set{20261125T005500}{4}% columns-from-a-sphere
\chapters@set{20261125T000000}{5}% slope-method
\chapters@set{20261126T030100}{6}% a-twist-for-each-prime
\chapters@set{20261125T003000}{7}% the-chain-and-the-square-root
\chapters@set{20261125T010000}{8}% rank-lemma-and-zero-estimate
\chapters@set{20261125T010500}{9}% poincare-series-and-near-orthogonality
\chapters@set{20261124T223000}{10}% elliptic-curves-for-the-proof
\chapters@set{20261125T009500}{11}% fock-model-and-theta-intertwiners
\chapters@set{20261124T223100}{12}% complex-multiplication-and-inert-reduction
\chapters@set{20261125T008000}{13}% one-cm-determinant
\chapters@set{20261124T214500}{14}% supersingular-theta-divisibility
\chapters@set{20261124T214600}{15}% the-isotropic-plane
\chapters@set{20261125T013700}{16}% pairings-on-a-formal-group-of-height-two
\chapters@set{20261125T013500}{17}% hecke-nondegeneracy-on-the-tree
\chapters@set{20261126T030200}{18}% the-bound-is-attained
\chapters@set{20261125T009000}{19}% kronecker-limit-and-chowla-selberg
\chapters@set{20261125T002000}{20}% definite-theta-lifts
\chapters@set{20261125T015500}{21}% siegel-weil-in-the-smallest-case
\chapters@set{20261125T002100}{22}% definite-theta-lifts-and-norms
\chapters@set{20261125T016500}{23}% hecke-operators-and-theta-lifts
\chapters@set{20261125T016700}{24}% divisibility-along-the-chain
\chapters@set{20261125T004000}{25}% denominators-and-sha
\chapters@set{20261125T004100}{26}% tate-twists-and-sha
\chapters@set{20261126T020430}{27}% the-minus-part-by-gauss-sums
\chapters@set{20261125T014100}{28}% congruence-modules
\chapters@set{20261125T014500}{29}% eisenstein-congruences-and-selmer-classes
\chapters@set{20261125T014700}{30}% pseudocharacters-and-the-length-bound
\chapters@set{20261125T011000}{31}% galois-representations-black-box
\chapters@set{20261126T030300}{32}% the-data
\chapters@set{20261126T020210}{33}% bessel-poincare-series-and-their-gram-operators
\chapters@set{20261126T030400}{34}% the-seesaw-identity
\chapters@set{20261126T030500}{35}% divisibility-of-the-composite-lift
\chapters@set{20261126T030600}{36}% theorem-a-short-forms
\chapters@set{20261126T030700}{37}% the-determinant-of-a-chain
\chapters@set{20261126T030800}{38}% the-cost-of-one-lift
\chapters@set{20261126T030900}{39}% denominators-of-projections-to-lifts
\chapters@set{20261126T031000}{40}% twists-with-few-congruences
\chapters@set{20261126T031100}{41}% theorem-c-and-the-main-inequality
\providecommand\chapnumber[1]{%
  \@ifundefined{chapters@#1}%
    {\PackageError{chapters}{Sheet #1 is not listed in the index}%
      {List it in 20261125T012000--grh-workbook__guide.org\MessageBreak
       and run scripts/make-chapters.py.}}%
    {\edef\thechap{\@nameuse{chapters@#1}}}}
\makeatother
\chapnumber{<sheet ID>}
\makeatletter
\providecommand\chapters@set[2]{\@namedef{chapters@#1}{#2}}
\chapters@set{20261125T001000}{1}% zeros-to-ramp-sums
\chapters@set{20261125T001100}{2}% inert-primes-and-estimate-b
\chapters@set{20261125T005000}{3}% galois-fourier-and-averaging
\chapters@set{20261125T005500}{4}% columns-from-a-sphere
\chapters@set{20261125T000000}{5}% slope-method
\chapters@set{20261126T030100}{6}% a-twist-for-each-prime
\chapters@set{20261125T003000}{7}% the-chain-and-the-square-root
\chapters@set{20261125T010000}{8}% rank-lemma-and-zero-estimate
\chapters@set{20261125T010500}{9}% poincare-series-and-near-orthogonality
\chapters@set{20261124T223000}{10}% elliptic-curves-for-the-proof
\chapters@set{20261125T009500}{11}% fock-model-and-theta-intertwiners
\chapters@set{20261124T223100}{12}% complex-multiplication-and-inert-reduction
\chapters@set{20261125T008000}{13}% one-cm-determinant
\chapters@set{20261124T214500}{14}% supersingular-theta-divisibility
\chapters@set{20261124T214600}{15}% the-isotropic-plane
\chapters@set{20261125T013700}{16}% pairings-on-a-formal-group-of-height-two
\chapters@set{20261125T013500}{17}% hecke-nondegeneracy-on-the-tree
\chapters@set{20261126T030200}{18}% the-bound-is-attained
\chapters@set{20261125T009000}{19}% kronecker-limit-and-chowla-selberg
\chapters@set{20261125T002000}{20}% definite-theta-lifts
\chapters@set{20261125T015500}{21}% siegel-weil-in-the-smallest-case
\chapters@set{20261125T002100}{22}% definite-theta-lifts-and-norms
\chapters@set{20261125T016500}{23}% hecke-operators-and-theta-lifts
\chapters@set{20261125T016700}{24}% divisibility-along-the-chain
\chapters@set{20261125T004000}{25}% denominators-and-sha
\chapters@set{20261125T004100}{26}% tate-twists-and-sha
\chapters@set{20261126T020430}{27}% the-minus-part-by-gauss-sums
\chapters@set{20261125T014100}{28}% congruence-modules
\chapters@set{20261125T014500}{29}% eisenstein-congruences-and-selmer-classes
\chapters@set{20261125T014700}{30}% pseudocharacters-and-the-length-bound
\chapters@set{20261125T011000}{31}% galois-representations-black-box
\chapters@set{20261126T030300}{32}% the-data
\chapters@set{20261126T020210}{33}% bessel-poincare-series-and-their-gram-operators
\chapters@set{20261126T030400}{34}% the-seesaw-identity
\chapters@set{20261126T030500}{35}% divisibility-of-the-composite-lift
\chapters@set{20261126T030600}{36}% theorem-a-short-forms
\chapters@set{20261126T030700}{37}% the-determinant-of-a-chain
\chapters@set{20261126T030800}{38}% the-cost-of-one-lift
\chapters@set{20261126T030900}{39}% denominators-of-projections-to-lifts
\chapters@set{20261126T031000}{40}% twists-with-few-congruences
\chapters@set{20261126T031100}{41}% theorem-c-and-the-main-inequality
\providecommand\chapnumber[1]{%
  \@ifundefined{chapters@#1}%
    {\PackageError{chapters}{Sheet #1 is not listed in the index}%
      {List it in 20261125T012000--grh-workbook__guide.org\MessageBreak
       and run scripts/make-chapters.py.}}%
    {\edef\thechap{\@nameuse{chapters@#1}}}}
\makeatother
\chapnumber{20261125T021600}
\renewcommand{\thesection}{\thechap.\arabic{section}}

\theoremstyle{definition}
\newtheorem{problem}{Problem}[section]
\theoremstyle{plain}
\newtheorem{theorem}{Theorem}[section]

\newcommand{\Lc}{\mathcal{L}}
\newcommand{\Bc}{\mathcal{B}}
\newcommand{\Wc}{\mathcal{W}}
\newcommand{\mumax}{\mu_{\max}}
\newcommand{\mumin}{\mu_{\min}}
\newcommand{\tl}[1]{\texttt{#1}}
\renewcommand{\Re}{\operatorname{Re}}

\setlength{\emergencystretch}{2em}

\title{The accounting in the Lean formalization}
\author{}
\date{November 25, 2026}

\begin{document}
\maketitle
\refstepcounter{chap}\label{chap}

\paragraph{Prerequisites.}
Chapter~\ref{as-chap}, all of it, and its notation.  From Chapter~\ref{zr-chap}: the functions \(G\) and \(G_\chi\) of~\eqref{zr-eq:G} and~\eqref{zr-eq:Gchi}, and Theorems~\ref{zr-thm:intervals} and~\ref{zr-thm:dirichlet-intervals}.  From Chapter~\ref{gf-chap}: Problems~\ref{gf-pr:thin} and~\ref{gf-pr:thick-move}.  From Chapter~\ref{sb-chap}: Theorem~\ref{sb-thm:near} and Problem~\ref{sb-pr:words}.  From Chapter~\ref{gl-chap}: Section~\ref{gl-sec:record} and Theorem~\ref{gl-thm:linkG}.  Chebyshev's bound \(\sum_{n \le y}\Lambda(n) \le 3y\) for \(y \ge 0\), as in Chapter~\ref{as-chap}.

\medskip
Chapter~\ref{as-chap} stated the inputs at one scale in a form close to \cite{Adr}, and proved the final theorem from them in the form the Lean formalization \cite{Lean} proves it.  What exactly does the formalization assume at one scale, and in what order does it make its choices?  This chapter records its accounting, and the problems check that the record gives the statements of Chapter~\ref{as-chap}.  The record has been read in the formalization's own statements; where we have read a statement but not its proof, we say so.  In this chapter \(c \ge 1\) is a constant of the record, which may depend on \(\varepsilon\) but not on \(x\), and \(s := \langle\mu(N) - h\rangle\).

\section{The record at one scale}\label{sec:record}

The formalization's statement at a scale \(x\) is one record for all translates \(\tau \in H\) at once (\tl{SizeData}\(_2\)).  For each \(\tau\) it has the words \(u\) and their numbers \(d_u\), \(T_u\), \(\mu(N_u)\), \(\mumax(N_u)\), \(h_u\), \(\mu(M_u)\), \(\mumax(M_u)\) and \(E_u\), and the numbers \(S\), \(\Gamma\), \(\Lc\), \(\mu(R^\vee)\) and \(\mumin(R^\vee)\) of Chapter~\ref{as-chap}.  It also has the excess colength \(I\) of the Hecke images, and the parameters \(\delta_F, \eta_F, e_F\) of the frame theorem.  Its inequalities are those of Chapter~\ref{as-chap}, inputs~\eqref{as-A1}--\eqref{as-A7} and~\eqref{as-A2p}--\eqref{as-A6v}, with the errors written in a few quantities: the precision \(L_a \le 2y\) of the linear term of the norms, the number \(X''\) and the weight \(b'\) of Chapter~\ref{gl-chap}, and a number \(L\) of the record.
\begin{itemize}
\item \eqref{as-A1}: \(S - \Gamma \le cN_Tmb\,x^{e_L\varepsilon}\), and also \(I \le cN_Tmb\,x^{e_L\varepsilon}\).
\item \eqref{as-A2}: the error depends on the word, \(e_2(u) = c\bigl(I + N_Tm(a/k + b')\log T_u + X''\bigr) + 1\); and \(e_2^+ = cmB_0\).
\item \eqref{as-A3}, in the form the frame theorem gives: if \(s - \mumin(R^\vee) \le Z\), then
\[
  s - \log\lvert\Wc\rvert - (1 - \delta_F)^{-1}\bigl(4u_F\log(1 + Z/u_F) + (\eta_F + 4\sqrt{e_F})Z\bigr) \le \mu(R^\vee) ,
\]
and the record has such a \(Z\) with \(0 < Z \le x^{C_H}b\), and \(0 \le \delta_F \le \bar\delta_F < 1\), \(0 \le \eta_F \le \bar\eta_F\), \(0 \le e_F \le \bar e_F \le \frac14\).
\item \eqref{as-A4}: \(e_4 = cN_TaL_a + cm(B_0 + N_c)x^{e_D\varepsilon}\).
\item \eqref{as-A5}, \eqref{as-A5m}: \(\lvert E_u - mAg\rvert \le cmB_0\).
\item \eqref{as-A6}: \(e_6 = c(L_a + m^2L_a/k)\); and \(\log T_u \le cmL\), \(e_T = c\log k\).
\item \eqref{as-A7}, \eqref{as-A7p}: \(e_7 = cN_TaL_a(1 + m/k)\) and \(e_7^+ = cN_TaL_a\).
\item \eqref{as-H}: \(e_H = c(N_TaL_a + mB_0)\).
\end{itemize}
The conclusion, for every translate, is the bound of Theorem~\ref{as-thm:acct} with an explicit error (\tl{each\_unit\_bound\_statement}).  Its proof is the three cancellations of Chapter~\ref{as-chap}, with the spread bounded as in Problem~\ref{as-pr:spread} (\tl{u\_function\_le\_mean}).

\begin{problem}\label{pr:key}
Let \(\ell_u\) be the right side of~\eqref{as-A6}, with \(e_6 = c(L_a + m^2L_a/k)\).  Then
\[
  \sum_t m\bigl(\varphi(a_t) - a_t\log\Omega\bigr) - A\ell_u + mAg = \sum_{j,t} M_{w_j}(a_t) + Ac\bigl(L_a + m^2L_a/k\bigr) .
\]
\end{problem}

\begin{problem}\label{pr:word-mean}
If the error \(e_2\) in~\eqref{as-A2} depends on the word, Theorem~\ref{as-thm:acct} and Problem~\ref{as-pr:spread} hold with \(e_2\) replaced by \(\langle e_2\rangle\).
\end{problem}

\begin{problem}\label{pr:Lo}
With the errors of the record, \(\langle e_2\rangle \le c\bigl(cN_Tmb\,x^{e_L\varepsilon} + cN_Tm^2(a/k + b')L + X''\bigr) + 1\).
\end{problem}

So the upper bound for \(\log T_u\) is used only to bound the error of~\eqref{as-A2}.

\begin{problem}\label{pr:phi}
For \(Z > 0\) the function \(v \mapsto v\log(1 + Z/v)\) is increasing on \((0, \infty)\).  Hence, if \(0 < u_F \le U\) and \(0 < Z \le x^{C_H}b\), then \(u_F\log(1 + Z/u_F) \le (U + b)\log(1 + x^{C_H})\).
\end{problem}

\begin{problem}\label{pr:e3-raw}
Suppose \(0 < u_F \le U\), \(\bar\delta_F \le \frac12\) and \((\bar\eta_F + 4\sqrt{\bar e_F})x^{C_H} \le 6\).  Then the form of~\eqref{as-A3} in the record implies~\eqref{as-A3} with \(e_3 = \log\lvert\Wc\rvert + 8\log(1 + x^{C_H})(U + b) + 12b\).
\end{problem}

The formalization obtains \((\bar\eta_F + 4\sqrt{\bar e_F})x^{C_H} \le 6\) from the construction of the frames.  Why does it need~\eqref{as-H}?

\begin{problem}\label{pr:H-needed}
Let \(\Wc = \{u, u'\}\) with \(d_u = d_{u'}\), let all the slopes \(\mu(N)\), \(\mumax(N)\), \(\mu(M)\), \(\mumax(M)\) be \(0\), let \(T_u = T_{u'}\), \(E_u = E_{u'} = mAg = \Gamma + \Lc\), and \(h_u = 0\), \(h_{u'} = D \ge 0\).  Then~\eqref{as-A2}, \eqref{as-A5}, \eqref{as-A7}, \eqref{as-A2p}, \eqref{as-A5m}, \eqref{as-A7p} and~\eqref{as-A6v} hold with all errors \(0\), and \(u_F = D/2\).
\end{problem}

So without~\eqref{as-H} the spread is not bounded by the other inputs: \(e_H\) must be at least \(D\) here.

\section{Where \texorpdfstring{\(\Gamma\)}{Gamma} and the slopes of the sources come from}\label{sec:sources}

At an inert place \(w\) the source lattice is the Hecke image twisted by the gains, \(N_{u,w} = \mathfrak p_w^{-\lfloor e_w\mathrm{sg}_u(w)\rfloor}N\!P_w\) (Section~\ref{gl-sec:record}).  In the formalization \(\Gamma\) is a weighted sum over the inert places of \(\sum_j\mathrm{gain}_j(p) - \tau(p)\): the colength of the Hecke image is paid inside \(\Gamma\), and appears in both~\eqref{as-A1} and~\eqref{as-A2} through \(\Gamma\).  Inputs~\eqref{as-A2} and~\eqref{as-A2p} then follow by linear algebra from three statements.
\begin{itemize}
\item The colength bound away from \(T_u\), Theorem~\ref{gl-thm:linkG}(iii); its excess over \(\tau\) is bounded by a multiple of \(I\).
\item The slope of \(N_u\) differs from that of its Hecke image, twisted by the gains, by at most a multiple of \(X''\).  The error is the rounding of the exponents, and the normalization at \(T_u\), at the level primes and at \(d\).  The formalization proves this from the divisibility branch (\tl{linkMu\_of\_linkDIV}).
\item At \(T_u\): the sublattice \(M_u^H\) of \(M_u\) at level \(T_u\) satisfies, exactly,
\[
  \mu(M_u^H) = \mu(M_u) - \frac{f\log(T_u - 1)}{T_u - 2} , \qquad \operatorname{rk} M_u^H = \Bigl(1 - \frac1{T_u - 1}\Bigr)^f\operatorname{rk} M_u ,
\]
and \(\mumax(M_u^H) \le \mumax(M_u)\) (\tl{adm26}).  The places over \(T_u\) cost at most \(cN_Tm(a/k + b')\log T_u\) per unit of rank.  We have read this statement, not its proof.
\end{itemize}

\begin{problem}\label{pr:rank-ratio}
If \(T \ge 3\) and \(d' = (1 - 1/(T - 1))^f d\), then \(0 \le 1 - d'/d \le f/(T - 1)\).
\end{problem}

\begin{problem}\label{pr:level-cost}
For every integer \(T \ge 3\), \(\dfrac{f\log(T - 1)}{T - 2} \le \dfrac{2f\log T}{T}\).
\end{problem}

So passing to level \(T_u\) lowers the slope by at most \(2f\log T_u/T_u\) and the rank by a proportion at most \(f/(T_u - 1)\); both are small, since \(T_u\) is large by~\eqref{as-A6}.

\section{The exponents}\label{sec:exponents}

The record writes the powers of \(x\) in two of its errors with exponents \(e_L\varepsilon\) and \(e_D\varepsilon\): in~\eqref{as-A1} on \(N_Tmb\), and in~\eqref{as-A4} on \(m(B_0 + N_c)\).  The formalization proves, for every translate, the bound of Theorem~\ref{as-thm:M}: the error is at most \(x^{(e_D + e_L + 3)\varepsilon}\Bc\) for large \(x\).  We have read the statement of the step that collects the errors into \(\Bc\), not its proof.  It then applies Theorem~\ref{as-thm:size}, which it proves with the constants of \cite[Lemma~6.1]{Zh} (\tl{proposition\_size\_bound}), with \(C = e_D + e_L + 3\).

The formalization works with a level: for a primitive character \(\psi\) of order \(o\), the primitive characters \(\psi'\) inducing the elements of the group generated by \(\psi\).  Its function is
\[
  G_m(y) := G(y) + \frac{1}{3o + 1}\sum_{\psi' \ne 1}\Re G_{\psi'}(y) ,
\]
and for \(\psi = \chi_8\), \(o = 2\), this is the function \(G + \frac17G_{\chi_8}\) of Chapter~\ref{as-chap}.

\begin{problem}\label{pr:Gm-positive}
For \(n \le y\), the coefficient of \(\Lambda(n)(y/n - 1)\) in \(G_m(y)\) is real and at least \(2(o + 1)/(3o + 1) \ge \frac23\).
\end{problem}

After averaging over the subgroup \(A_0\), the characters not trivial on \(A_0\) contribute zero, by orthogonality, and \(4o\) characters remain.  Those that induce the level give \(G_m\), and the formalization bounds \(G_m(x)\) by \(x^{1/2+(C+3)\varepsilon}\) plus the sum over the others, which Estimate~B bounds by \(4o\,x^{1/2+\varepsilon}\).

\begin{problem}\label{pr:absorb}
Let \(C \ge 0\), \(o \ge 1\), \(x \ge 1\), \(\varepsilon > 0\) and \(1 + 4o \le x^\varepsilon\).  Then \(x^{1/2+(C+3)\varepsilon} + 4o\,x^{1/2+\varepsilon} \le x^{1/2+(C+4)\varepsilon}\).
\end{problem}

With \(C = e_D + e_L + 3\), \(e_D = 2\) and \(e_L = 0\), this is the formalization's bound
\[
  G_m(x) \le x^{1/2 + (e_D + e_L + 7)\varepsilon} = x^{1/2+9\varepsilon}
\]
for every large \(x\) in the set of scales of the family; this is \tl{Chebyshev\allowbreak{}Inequality\allowbreak{}On\allowbreak{}Windows}.

\section{The costs of the averaging}\label{sec:costs}

The formalization does not use Problem~\ref{gf-pr:thick-move}.  It moves the thick weights to \(b\) with a Lipschitz bound in the weight, and states the costs as condition (V3) of Chapter~\ref{as-chap} (\tl{CostWindows}).

\begin{problem}\label{pr:lip}
For \(A > 0\) and \(w, w' > 0\),
\(
  \bigl\lvert wG(A/w) - w'G(A/w')\bigr\rvert \le 3A\lvert w - w'\rvert/\min(w, w')
\).
The same holds with \(G_\chi\) in place of \(G\), for \(\chi\) primitive and nonprincipal.
\end{problem}

With Problem~\ref{gf-pr:thin}, the formalization bounds the cost of one pair \((t, j)\) by \(2f(x+1)b\vartheta + 6f(x+1)(\Delta + 1)\), which (V3) makes at most \(2b/x\).  It moves \(a_t/b\) to \(x\) with a bound for the derivative of \(G_\chi\) that involves \(\lvert L'/L(1, \chi)\rvert\), and takes \(\lvert L'/L(1, \chi)\rvert \le C_L\log x\) as a condition on the data at \(x\).  In the thin-scale bound, \(\vartheta = \exp(-c\,x^{\varepsilon/8})\).

\section{The choices of \texorpdfstring{\(m\), \(k\) and \(n\)}{m, k and n}}\label{sec:choices}

The normalization needs only the inequalities~\eqref{as-eq:S}.  Where do the sizes \(k \approx x^{9/2}\) and \(n \approx x^2\) come from?  Let \(\delta_1 := x^{-\varepsilon/8}\) as in Chapter~\ref{as-chap}.

\begin{problem}\label{pr:compatible}
A real number \(k\) with \(xm/k \le 1\) and \(k/n^2 \le x^{1/2-\varepsilon}\) exists if and only if \(n^2 \ge m\,x^{1/2+\varepsilon}\).  With \(m \ge x^{1/2-\varepsilon/16}\) this forces \(n \ge x^{1/2}\).
\end{problem}

Since \(k\) is a dimension it must also be an integer; at the scale of the argument the interval \([xm, n^2x^{1/2-\varepsilon}]\) is long, so this costs nothing.

The lower bound for \(k\) comes from the rank hypothesis.  The sources must have total rank at least \(x^{C_H}\) times the dimension of the target, each has rank at most \(C_0\) times that dimension divided by \(\Gamma_0^2 = \exp(4C_\Gamma m^5ny)\), the constant of near orthogonality (Theorem~\ref{sb-thm:near}), and there are at most \(2^{k-3}\) of them; so \(k \ge m^5n\) for large \(x\) (Problem~\ref{sb-pr:words}).  This bound is needed by~\eqref{as-A3}, not by the normalization.

\begin{problem}\label{pr:x2}
If \(k \ge m^5n\), \(k \le n^2x^{1/2}\) and \(m \ge x^{1/2-\varepsilon/16}\), then \(n \ge x^{2-5\varepsilon/16}\) and \(k \ge x^{9/2-5\varepsilon/8}\).
\end{problem}

The choices of \cite[Prop.~5.15]{Adr} and \cite[\S I.5]{Rub} are \(m \in [\sqrt{x\delta_1}, 2\sqrt{x\delta_1}]\), \(k \in [x^{9/2+10\varepsilon}, 2x^{9/2+10\varepsilon}]\) and \(n \in [x^{2+8\varepsilon}, 2x^{2+8\varepsilon}]\), with \(n\) a power of \(2\) times a bounded integer.

\begin{problem}\label{pr:substitute}
With these choices, \(x/m\), \(xm/k\), \(m\) and \(k/n^2\) are at most \(x^{1/2+\varepsilon/16}\), \(2x^{-3-161\varepsilon/16}\), \(2x^{1/2-\varepsilon/16}\) and \(2x^{1/2-6\varepsilon}\); for \(\varepsilon = \frac1{100}\) the exponents are \(0.500625\), \(-3.100625\), \(0.499375\) and \(0.44\).  Moreover~\eqref{as-eq:S} holds for large \(x\) if \(\delta \ge x^{-\varepsilon}\), \(N_T \ge \frac12 x^{-\varepsilon/8}n\) and \(N_2 \le n\).
\end{problem}

\begin{problem}\label{pr:139}
With these choices, \(4C_\Gamma m^5ny \le 256\,C_\Gamma\,x^{9/2+139\varepsilon/16}\).  Since \(\frac{139}{16} < 10\), the inequality of Problem~\ref{sb-pr:words} holds for large \(x\).
\end{problem}

In the formalization these are the choices made when the family is built (\tl{RankVal}).

\section{The order of choices}\label{sec:order}

The formalization proves the following, for each primitive character \(\psi\) (proved by \tl{topDesWin}).  There are \(e_D\), \(e_L\) and \(\varepsilon_1 > 0\) such that for every \(\varepsilon \in (0, \varepsilon_1)\) there is a family of data whose set of scales contains \([Y, Y^{3/2}]\) for arbitrarily large \(Y\), and constants and \(x_0\), such that the record of Section~\ref{sec:record} exists at every \(x \ge x_0\) in that set.  It takes \(e_D = 2\) and \(e_L = 0\).  In the proof the constants are chosen before the family, which is large relative to them.  The record at \(x\) ties together one member of the family for each translate.  By contrast, \cite[Hyp.~18.10]{Adr} asks for the inputs at every admissible datum.

\begin{problem}\label{pr:small-eps}
Let \(c_0 \in \mathbb R\), \(\delta_0 > 0\) and \(\varepsilon_1 > 0\).  Then \(\varepsilon := \min\bigl(\varepsilon_1/2, \delta_0/(\lvert c_0\rvert + 1)\bigr)\) satisfies \(0 < \varepsilon < \varepsilon_1\) and \(c_0\varepsilon \le \delta_0\).
\end{problem}

\begin{problem}\label{pr:zero-free}
Fix \(e_D, e_L \ge 0\) and \(\varepsilon_1 > 0\), and suppose that for every \(\varepsilon \in (0, \varepsilon_1)\) there are arbitrarily large \(Y\) with \(G(x) \le x^{1/2+(e_D+e_L+7)\varepsilon}\) for all \(x \in [Y, Y^{3/2}]\).  Let \(\Re s > \frac12\) and \(\theta := \frac12 + \frac12\min(\Re s - \frac12, \frac12)\).  Then there are arbitrarily large \(Y\) with \(G(x) \le x^\theta\) on \([Y, Y^{3/2}]\), and \(\theta < \Re s\); so \(\zeta(s) \ne 0\).
\end{problem}

This is the formalization's last step (\tl{level\_of\_tiedWin}), for \(G_m\) and Theorem~\ref{zr-thm:dirichlet-intervals} in place of \(G\) and Theorem~\ref{zr-thm:intervals}.  The choice of \(\varepsilon\) depends on \(s\) only through \(e_D + e_L + 7\), which is why \(e_D\) and \(e_L\) come before \(\varepsilon\).

\section*{Hints}

\subsection*{\S\ref{sec:record}}
\begin{description}
\item[\ref{pr:key}] Expand \(\varphi\) and \(\log\Omega\), and use \(\lvert\Phi\rvert = f\), as in Problem~\ref{as-pr:levels}.
\item[\ref{pr:word-mean}] In Problem~\ref{as-pr:A127} keep \(e_2(u)\), and take the weighted mean.
\item[\ref{pr:Lo}] Problem~\ref{pr:word-mean}, the bound for \(I\), and \(\log T_u \le cmL\).
\item[\ref{pr:phi}] The derivative is \(\log(1 + t) - t/(1 + t)\) with \(t = Z/v\).  Then compare with \(v = U + b\).
\item[\ref{pr:e3-raw}] Problem~\ref{pr:phi}, \((1 - \delta_F)^{-1} \le 2\) and \(\sqrt{e_F} \le \sqrt{\bar e_F}\).
\item[\ref{pr:H-needed}] \(\max_u(\mumax(N_u) - h_u) = 0\).
\end{description}

\subsection*{\S\ref{sec:sources}}
\begin{description}
\item[\ref{pr:rank-ratio}] Bernoulli's inequality.
\item[\ref{pr:level-cost}] For \(T \ge 4\), \(T \le 2(T - 2)\); check \(T = 3\) directly.
\end{description}

\subsection*{\S\ref{sec:exponents}}
\begin{description}
\item[\ref{pr:Gm-positive}] There are \(o - 1\) characters \(\psi' \ne 1\), and \(\lvert\psi'(n)\rvert \le 1\).
\item[\ref{pr:absorb}] \(x^{1/2+\varepsilon} \le x^{1/2+(C+3)\varepsilon}\).
\end{description}

\subsection*{\S\ref{sec:costs}}
\begin{description}
\item[\ref{pr:lip}] The function \(w \mapsto wG(A/w)\) is continuous, and away from the points \(A/n\) its derivative is \(A/w - \sum_{n \le A/w}\Lambda(n)\).
\end{description}

\subsection*{\S\ref{sec:choices}}
\begin{description}
\item[\ref{pr:x2}] Eliminate \(k\).
\item[\ref{pr:139}] \(m \le 2x^{1/2-\varepsilon/16}\) and \(n \le 2x^{2+8\varepsilon}\).
\end{description}

\subsection*{\S\ref{sec:order}}
\begin{description}
\item[\ref{pr:small-eps}] \(c_0 \le \lvert c_0\rvert\).
\item[\ref{pr:zero-free}] Problem~\ref{pr:small-eps} with \(c_0 = e_D + e_L + 7\), and Theorem~\ref{zr-thm:intervals}.
\end{description}

\begin{thebibliography}{Lean}
\bibitem[Adr]{Adr} J.~Adriazola, \emph{The Alp\"oge proof of the Riemann hypothesis}, manuscript, revised 24 November 2026, \url{https://github.com/jimjamzola/RH-Alpoge-Reconstruction}.
\bibitem[Lean]{Lean} The Lean community, \emph{riemann-hypothesis}, a Lean formalization of Alp\"oge's argument, \url{https://github.com/leanprover-community/riemann-hypothesis}.
\bibitem[Rub]{Rub} \emph{The Alp\"oge--Claude proof of the Riemann hypothesis explained for the Riemann zeta function}, written by AI models at the behest of M.~O.~Rubinstein, working draft, 21 November 2026, \url{https://www.math.uwaterloo.ca/~mrubinst/alpoge_claude_grh/zeta-paper.pdf}.
\bibitem[Zh]{Zh} R.~Zhang, \emph{A preliminary note on Alp\"oge's proposed proof of the generalized Riemann hypothesis}, preprint, 2026, hexagon:2611.00809.
\end{thebibliography}

\end{document}
